\heading{经济学原理（II）-模拟题4-期中}
\noteinfo{这套题继续提高难度，强调“一个题里同时考多个知识点”。尤其是第 17--21 题，建议把题目先翻译成经济学语言，再动手计算。}

\instructions{建议答题时间 120 分钟。图形题未画图也可作答，但必须用文字清楚描述哪条曲线移动、为什么移动，以及均衡变量如何变化。}

\textbf{一、单项选择题}

\begin{question}
若某政策同时提高了当前总产出水平并减少了贫困率，但基尼系数几乎不变，则最合理的理解是
\begin{choiceoptions}[1]
  \optionitem{该政策必然没有改善低收入者福利}
  \optionitem{贫困与不平等衡量的是完全相同的问题}
  \optionitem{政策可能抬高了所有收入水平，尤其使部分贫困者越过贫困线，但对整体相对分配影响有限}
  \optionitem{基尼系数必然算错了}
\end{choiceoptions}
\begin{solution}
贫困关注是否低于某门槛，不平等关注相对分配。政策可以显著降低贫困，却对整体相对分配结构影响不大。\anstext{C}
\end{solution}
\end{question}

\begin{question}
下列哪一项最可能\textbf{计入 GDP 平减指数，但不直接计入 CPI}？
\begin{choiceoptions}[2]
  \optionitem{居民购买的进口手机价格上涨}
  \optionitem{政府采购的国产军用设备价格上涨}
  \optionitem{居民租房价格上涨}
  \optionitem{进口原油零售价上涨并被居民直接消费}
\end{choiceoptions}
\begin{solution}
GDP 平减指数覆盖国内生产的政府购买品，而 CPI 只覆盖居民消费篮子，因此政府采购国产军用设备价格上涨会进入平减指数而不进入 CPI。\anstext{B}
\end{solution}
\end{question}

\begin{question}
若某国资本份额为 $\alpha=0.4$，在人均 Cobb--Douglas 生产函数下，黄金规则储蓄率通常是
\begin{choiceoptions}[2]
  \optionitem{0.2}
  \optionitem{0.4}
  \optionitem{0.6}
  \optionitem{与 $\alpha$ 无关}
\end{choiceoptions}
\begin{solution}
在 $y=k^{\alpha}$ 下，黄金规则储蓄率为 $\alpha$。\anstext{B}
\end{solution}
\end{question}

\begin{question}
若政府从预算盈余转为预算赤字，而私人储蓄和投资需求曲线都不变，则
\begin{choiceoptions}[1]
  \optionitem{国民储蓄减少，真实利率上升，投资下降}
  \optionitem{国民储蓄增加，真实利率下降，投资上升}
  \optionitem{只有名义利率变化，真实利率不变}
  \optionitem{只有公共储蓄变化，投资一定不受影响}
\end{choiceoptions}
\begin{solution}
预算赤字减少公共储蓄，从而减少国民储蓄，使供给左移、真实利率上升、投资下降。\anstext{A}
\end{solution}
\end{question}

\begin{question}
关于链式加权真实 GDP，下列说法最准确的是
\begin{choiceoptions}[1]
  \optionitem{它完全放弃价格信息，只追踪数量}
  \optionitem{它试图减少固定基期价格在长期比较中变得不现实的问题}
  \optionitem{它等同于 CPI 的计算方法}
  \optionitem{它保证真实 GDP 永远高于名义 GDP}
\end{choiceoptions}
\begin{solution}
链式加权的目的是动态更新权重，减少固定基期价格长期失真的问题。\anstext{B}
\end{solution}
\end{question}

\begin{question}
若一只股票的 $\beta=1.5$，最恰当的解释是
\begin{choiceoptions}[2]
  \optionitem{该股票没有非系统性风险}
  \optionitem{市场上涨 1\% 时，该股票一定上涨 1.5\%}
  \optionitem{该股票对市场共同波动更敏感，系统性风险暴露较高}
  \optionitem{该股票的标准差一定是市场的 1.5 倍}
\end{choiceoptions}
\begin{solution}
$\beta$ 描述的是对市场共同波动的敏感度，不等于“每次都严格涨 1.5 倍”。\anstext{C}
\end{solution}
\end{question}

\begin{question}
关于有效市场假说与资产泡沫，下列哪项说法更稳妥？
\begin{choiceoptions}[1]
  \optionitem{只要存在泡沫，市场就必然完全无效}
  \optionitem{若市场半强有效，就不可能出现任何价格偏离}
  \optionitem{即便市场对信息反应很快，也仍可能因为预期、信念与套利约束而出现泡沫}
  \optionitem{有效市场假说说明所有投资者都完全理性且意见一致}
\end{choiceoptions}
\begin{solution}
市场对信息快速反映，并不自动排除泡沫；现实中可能存在套利限制、信念差异和预期自我实现。\anstext{C}
\end{solution}
\end{question}

\begin{question}
若某国实际 GDP 年增长率为 4\%，人口增长率为 1\%，则人均实际 GDP 年增长率近似为
\begin{choiceoptions}[2]
  \optionitem{1\%}
  \optionitem{3\%}
  \optionitem{4\%}
  \optionitem{5\%}
\end{choiceoptions}
\begin{solution}
近似公式：人均增长率 $\approx$ 总量增长率 $-$ 人口增长率 $=3\%$。\anstext{B}
\end{solution}
\end{question}

\begin{question}
若可贷资金市场上的真实利率下降，同时预期通胀率上升且幅度更大，则名义利率
\begin{choiceoptions}[2]
  \optionitem{一定下降}
  \optionitem{一定上升}
  \optionitem{可能上升也可能下降，取决于二者净效应}
  \optionitem{一定不变}
\end{choiceoptions}
\begin{solution}
名义利率近似为 $i\approx r+\pi^e$。真实利率下降但预期通胀上升更多，净效应取决于幅度比较。\anstext{C}
\end{solution}
\end{question}

\begin{question}
在两期消费模型中，边际消费倾向小于 1 的根本原因是
\begin{choiceoptions}[1]
  \optionitem{消费者总是更喜欢未来消费}
  \optionitem{消费者会把部分新增现期收入分配到未来，从而进行跨期平滑}
  \optionitem{现期收入永远低于未来收入}
  \optionitem{效用函数必须是二次函数}
\end{choiceoptions}
\begin{solution}
只要消费者能进行跨期选择，就会把新增收入的一部分留给未来使用，因此现期 MPC 通常小于 1。\anstext{B}
\end{solution}
\end{question}

\textbf{二、简答题}

\begin{question}
为什么“最高收入 5\% 的收入份额”这一类额外信息，会显著改善基尼系数的近似估计？
\begin{solution}
因为收入分布往往具有厚尾特征，最高收入组内部并不均匀。如果只知道五等份，最高 20\% 内部的巨大差异会被抹平，从而低估不平等。进一步知道最高 5\% 或 1\% 的份额，就能把尾部再细分一次，更真实地刻画高收入集中度，因此洛伦兹曲线近似和基尼系数估计都会更准确。\end{solution}
\end{question}

\begin{question}
请比较 GDP 与 GNP 在国际资本流动时代各自的优点。为什么一个国家分析“境内生产能力”与“居民所得”时，往往需要同时看二者？
\begin{solution}
GDP 反映境内生产规模与本土就业、税基、产业活动，适合分析“国家在本地生产了多少”；GNP 反映本国居民最终获得多少生产收入，适合分析“居民从全球生产中分到多少”。在国际资本流动时代，一个国家可能境内有大量外资生产，也可能本国居民在海外拥有大量资产，因此 GDP 与 GNP 会出现明显差异。要理解一国经济实力与居民福利，往往必须同时看二者。\end{solution}
\end{question}

\begin{question}
为什么黄金规则追求的是“最大稳态消费”而不是“最大稳态产出”？请用经济学语言解释其含义。
\begin{solution}
产出不是最终目的，消费才更直接对应居民福利。提高储蓄率会提升资本与产出，但也意味着当前和稳态下都要拿出更多资源做投资。如果储蓄率过高，虽然产出更大，但留给居民消费的部分不一定更多。黄金规则正是在“更多资本带来的额外产出”与“为了维持更多资本而牺牲的资源”之间求平衡，使长期稳态消费最大。\end{solution}
\end{question}

\begin{question}
“预算赤字的挤出效应”与“投资税收抵免刺激投资”看似方向相反。请说明二者可以怎样同时成立。
\begin{solution}
预算赤字通过减少国民储蓄，使可贷资金供给左移，抬高真实利率，从而压低投资；投资税收抵免则直接提高企业投资意愿，使投资需求右移。二者作用在不同曲线上，因此可以同时成立。结果上，真实利率通常会上升，但投资总量是上升还是下降，需要看哪一种力量更强。\end{solution}
\end{question}

\textbf{三、计算和分析题}

\begin{question}
某经济体有三家企业，企业 3 的产品全部卖给家庭。已知：
\begin{center}
\begin{tabular}{c|cccc}
 & 企业1 & 企业2 & 企业3 & 经济整体 \\\hline
销售收入 & 600 & 1000 & \blank[1.2cm] & -- \\
原料投入 & 0 & \blank[1.2cm] & 1000 & -- \\
劳动报酬 & 300 & 180 & 400 & \blank[1.2cm] \\
地租 & 60 & 50 & 80 & \blank[1.2cm] \\
利息 & 40 & 20 & 40 & \blank[1.2cm] \\
利润 & \blank[1.2cm] & 150 & 280 & \blank[1.2cm] \\
增加值 & \blank[1.2cm] & \blank[1.2cm] & \blank[1.2cm] & \blank[1.2cm]
\end{tabular}
\end{center}
请补全并求 GDP。
\begin{solution}
企业 1 利润为
\[
600-300-60-40=200,
\]
故增加值为 600。企业 2 原料投入来自企业 1，为 600，故增加值为
\[
1000-600=400.
\]
企业 3 增加值为
\[
400+80+40+280=800.
\]
其销售收入为
\[
1000+800=1800.
\]
经济整体劳动报酬 $=880$，地租 $=190$，利息 $=100$，利润 $=630$。GDP 为
\[
600+400+800=1800.
\]
\ansmath{GDP=1800}
\end{solution}
\end{question}

\begin{question}
某国五等份收入份额为：6\%、11\%、16\%、23\%、44\%。
\begin{enumerate}[label=(\alph*),leftmargin=2.5em]
  \item 写出洛伦兹曲线关键点；
  \item 计算基尼系数；
  \item 若政府实施一次只针对最低 20\% 的转移支付，使最低组份额升至 8\%，最高组份额降至 42\%，其余三组不变，则新的基尼系数是多少？
\end{enumerate}
\begin{solution}
初始累积收入份额为
\[
0,\ 0.06,\ 0.17,\ 0.33,\ 0.56,\ 1.
\]
因此点为
\[
(0,0),(0.2,0.06),(0.4,0.17),(0.6,0.33),(0.8,0.56),(1,1).
\]
面积：
\[
A_L=0.2\times\frac{0+0.06}{2}+0.2\times\frac{0.06+0.17}{2}+0.2\times\frac{0.17+0.33}{2}+0.2\times\frac{0.33+0.56}{2}+0.2\times\frac{0.56+1}{2}=0.322.
\]
故初始基尼系数为
\[
G=1-2\times0.322=0.356.
\]

政策后累积份额为
\[
0,\ 0.08,\ 0.19,\ 0.35,\ 0.58,\ 1,
\]
新面积
\[
A'_L=0.2\times\frac{0+0.08}{2}+0.2\times\frac{0.08+0.19}{2}+0.2\times\frac{0.19+0.35}{2}+0.2\times\frac{0.35+0.58}{2}+0.2\times\frac{0.58+1}{2}=0.34.
\]
故新基尼系数
\[
G'=1-2\times0.34=0.32.
\]
\ansmath{G=0.356,\quad G'=0.320}
\end{solution}
\end{question}

\begin{question}
一个经济体生产三种商品，数据如下：
\begin{center}
\begin{tabular}{c|cc|cc|cc}
 & \multicolumn{2}{c|}{2023年} & \multicolumn{2}{c|}{2024年} & \multicolumn{2}{c}{2025年} \\
商品 & 产量 & 价格 & 产量 & 价格 & 产量 & 价格 \\\hline
苹果 & 10 & 4 & 12 & 5 & 15 & 6 \\
梨 & 20 & 2 & 18 & 2 & 20 & 3 \\
葡萄 & 30 & 1 & 35 & 2 & 40 & 2
\end{tabular}
\end{center}
\begin{enumerate}[label=(\alph*),leftmargin=2.5em]
  \item 计算 2024 至 2025 年名义 GDP 增长率；
  \item 以 2024 年为基期，求 2023 与 2025 年真实 GDP；
  \item 分别求 2025 年的 GDP 平减指数与若固定消费篮子为苹果 1、梨 2、葡萄 3 时的 CPI（以 2024 年为基期）；
  \item 比较二者大小并解释原因。
\end{enumerate}
\begin{solution}
2024 年名义 GDP：
\[
12\times5+18\times2+35\times2=166.
\]
2025 年名义 GDP：
\[
15\times6+20\times3+40\times2=230.
\]
增长率：
\[
\frac{230-166}{166}\times100\%=38.55\%.
\]

以 2024 年价格算真实 GDP：
\[
RGDP_{2023|2024}=10\times5+20\times2+30\times2=150,
\]
\[
RGDP_{2025|2024}=15\times5+20\times2+40\times2=195.
\]

GDP 平减指数：
\[
P^{GDP}_{2025|2024}=\frac{230}{195}\times100\approx117.95.
\]
固定篮子成本：2024 年为
\[
1\times5+2\times2+3\times2=15,
\]
2025 年为
\[
1\times6+2\times3+3\times2=18,
\]
故
\[
CPI_{2025|2024}=\frac{18}{15}\times100=120.
\]
这里 CPI 略高于 GDP 平减指数，是因为固定消费篮子对涨价较多的商品赋予了不同权重，而平减指数使用的是当期产出权重。\ansmath{\text{名义 GDP 增长率}\approx 38.55\%,\ RGDP_{2023|2024}=150,\ RGDP_{2025|2024}=195}
\ansmath{P^{GDP}_{2025|2024}\approx117.95,\ CPI_{2025|2024}=120}
\end{solution}
\end{question}

\begin{question}
设某经济体的人均生产函数为
\[
y=k^{0.4},
\]
储蓄率 $s=0.2$，人口增长率 $n=0.03$，折旧率 $\delta=0.05$。回答：
\begin{enumerate}[label=(\alph*),leftmargin=2.5em]
  \item 求稳态条件与稳态资本密度 $k^*$；
  \item 求稳态人均消费 $c^*$；
  \item 求黄金规则储蓄率；
  \item 若政府把储蓄率从 0.2 提高到 0.4，稳态人均产出会怎样变化？稳态人均消费一定上升吗？
\end{enumerate}
\begin{solution}
稳态条件是
\[
0.2k^{0.4}=0.08k.
\]
即
\[
k^{0.6}=\frac{0.2}{0.08}=2.5,
\qquad
k^*=2.5^{\frac{1}{0.6}}=2.5^{5/3}\approx 4.604.
\]
稳态人均产出为
\[
y^*=(k^*)^{0.4}=2.5^{2/3}\approx 1.842.
\]
稳态人均消费为
\[
c^*=(1-s)y^*=0.8\times1.842\approx 1.474.
\]
黄金规则储蓄率在 Cobb--Douglas 下等于资本份额，即
\[
s_{gold}=0.4.
\]
若把储蓄率从 0.2 提高到 0.4，稳态人均产出一定上升，因为稳态资本密度更高；但稳态人均消费不总是一定上升，只有当原储蓄率低于黄金规则、上调后不超过最优附近时，消费才会上升。本题中 0.2 低于黄金规则，提高到 0.4 正好到达黄金规则，因此稳态消费也会上升到最大值。\ansmath{k^*\approx 4.604,\quad c^*\approx 1.474,\quad s_{gold}=0.4}
\end{solution}
\end{question}

\begin{question}
考虑可贷资金市场与费雪方程。起初均衡真实利率为 3\%，预期通胀率为 2\%。随后政府实行紧缩财政，使公共储蓄增加，均衡真实利率下降到 2\%；与此同时，央行成功把预期通胀率锚定到 1\%。
\begin{enumerate}[label=(\alph*),leftmargin=2.5em]
  \item 初始名义利率是多少？
  \item 政策后名义利率是多少？
  \item 请用“储蓄供给变化 + 费雪方程”完整解释名义利率为什么这样变化。
\end{enumerate}
\begin{solution}
初始名义利率约为
\[
i_0\approx r_0+\pi^e_0=3\%+2\%=5\%.
\]
政策后名义利率约为
\[
i_1\approx r_1+\pi^e_1=2\%+1\%=3\%.
\]
紧缩财政提高公共储蓄，使可贷资金供给右移，真实利率下降；同时预期通胀下降，进一步压低名义利率。于是名义利率从 5\% 降至 3\%。\ansmath{i_0\approx5\%,\quad i_1\approx3\%}
\end{solution}
\end{question}

\begin{question}
两种资产 $C$ 与 $D$ 的收益率如下，概率各为 0.5：
\begin{center}
\begin{tabular}{c|cc}
状态 & $C$ & $D$\\\hline
1 & 40\% & -10\%\\
2 & -20\% & 20\%
\end{tabular}
\end{center}
\begin{enumerate}[label=(\alph*),leftmargin=2.5em]
  \item 求两种资产的期望收益率；
  \item 若对 $C$ 的配置权重为 $w$，求使组合无风险的 $w$；
  \item 求此时的无风险收益率；
  \item 解释为什么这一无风险组合并不意味着现实世界中“金融市场不再需要风险溢价”。
\end{enumerate}
\begin{solution}
\[
E(r_C)=0.5\times40\%+0.5\times(-20\%)=10\%,
\]
\[
E(r_D)=0.5\times(-10\%)+0.5\times20\%=5\%.
\]
状态 1 的组合收益率为
\[
r_1=0.4w-0.1(1-w)=0.5w-0.1,
\]
状态 2 为
\[
r_2=-0.2w+0.2(1-w)=0.2-0.4w.
\]
令两者相等：
\[
0.5w-0.1=0.2-0.4w\Rightarrow 0.9w=0.3\Rightarrow w=\frac13.
\]
无风险收益率为
\[
r_f=0.5\times\frac13-0.1=\frac{1}{15}\approx 6.67\%.
\]
即使存在某一对资产能被构造成局部无风险组合，也不意味着所有资产和所有状态都可完美对冲；现实中相关性会变化，交易成本、卖空约束和系统性冲击依然存在，因此风险溢价并不会消失。\ansmath{E(r_C)=10\%,\quad E(r_D)=5\%,\quad w=\frac13,\quad r_f\approx 6.67\%}
\end{solution}
\end{question}

\begin{question}
两期消费模型中，消费者效用函数为
\[
u(c_1,c_2)=\sqrt{c_1}+\sqrt{c_2},
\]
初始资产为 0，真实利率为 $r$，并满足 $y_2=(1+r)y_1$。
\begin{enumerate}[label=(\alph*),leftmargin=2.5em]
  \item 写出欧拉方程；
  \item 写出终身预算约束；
  \item 解出 $c_1$ 与 MPC；
  \item 说明为什么这个例子里的 MPC 仍小于 1。
\end{enumerate}
\begin{solution}
欧拉方程：
\[
\frac{1}{2\sqrt{c_1}}=(1+r)\frac{1}{2\sqrt{c_2}}
\Rightarrow c_2=(1+r)^2c_1.
\]
预算约束：
\[
c_1+\frac{c_2}{1+r}=y_1+\frac{y_2}{1+r}=2y_1.
\]
代入得
\[
c_1+(1+r)c_1=(2+r)c_1=2y_1,
\]
故
\[
c_1=\frac{2}{2+r}y_1.
\]
于是
\[
MPC=\frac{\partial c_1}{\partial y_1}=\frac{2}{2+r}<1.
\]
它小于 1，是因为消费者不会把新增现期收入全部花在当期，而会把一部分留给未来消费，以平滑跨期边际效用。\ansmath{c_1=\frac{2}{2+r}y_1,\quad MPC=\frac{2}{2+r}}
\end{solution}
\end{question}

